\documentclass{situationist-preprint}

\title{Naïve Realism Restored}
\msid{MND-16-a}
\msfamily{MND-16}
\msdiscipline{Philosophy of Mind}
\mstype{Principal manuscript}
\mskeywords{perception; naïve realism; disjunctivism; argument from hallucination; common kind; screening off; seemings; indiscriminability}

\begin{document}
\maketitle

\begin{abstract}
Naïve realism holds that in veridical perception the subject is presented with mind-independent objects and their properties, which partly constitute the experience and its phenomenal character. The argument from hallucination is the main obstacle: since a hallucination can be indistinguishable from a perception and can have the same proximal cause, it is said, the two must be experiences of the same kind, and that kind cannot involve the object. We argue that naïve realism is true and that hallucinations are experiences of a different kind from veridical perceptions. The argument from hallucination relies on a principle of common kind, that the same proximal cause yields an experience of the same fundamental kind. We show that this principle is false for relational events, and that the indiscriminability of the two experiences is explained by what they share at another level: both produce in the subject a seeming that things are thus and so, an attitude common to perception and hallucination. Distinguishing experiences, which differ in kind, from seemings, which can be common, answers the screening-off objection and makes naïve realism compatible with an internalist epistemology of perceptual justification.
\end{abstract}

\section{The problem}

When one sees a lemon on a table, one seems to be confronted with the lemon itself: its color, shape and position are presented as there to be seen, and the experience seems to be constituted by them. \emph{Naïve realism} takes this appearance at face value. On the naïve realist view, a veridical perceptual experience is a relation between a subject and a mind-independent object, and the object and its properties partly constitute the experience and its phenomenal character \citep{Martin2002,Campbell2002,Brewer2011}. The view contrasts with sense-datum theories, on which one is directly aware only of mental items \citep{Russell1912,Ayer1940,Robinson1994}, and with representationalism, on which the experience represents the lemon as being a certain way, and could have the same character whether or not the lemon is there \citep{Tye1995,Siegel2010}.

The question whether perception is direct is the question whether naïve realism, or something like it, is true. We argue that it is. The classical obstacle is the argument from hallucination, and we answer it by giving an account of what hallucination is.

\section{A short history}

Hume observed that the table we see seems to diminish as we move away from it, while the real table does not change, and concluded that what is present to the mind is only an image (\emph{Enquiry} XII). The arguments from illusion and from hallucination became, in the hands of \citet{Russell1912} and \citet{Ayer1940}, the standard route to sense data. \citet{Moore1903} had already noted the elusiveness of the experience as distinct from its object, the observation that later became the transparency of experience.

The modern revival of direct realism began with the observation that the argument from hallucination assumes, without argument, that the experiences in the good and bad cases are of the same kind. \citet{Hinton1973} proposed that a statement of how things perceptually appear to a subject is equivalent to a disjunction: either she perceives that things are thus, or it is merely as if she does. \citet{Snowdon1981} applied the disjunctive analysis to the causal theory of perception, and \citet{McDowell1982} to the epistemology of perception, arguing that in the good case the subject's experience makes a fact manifest, which a hallucination does not. \citet{Martin2002,Martin2004,Martin2006} developed the view as a thesis about the nature of experience, and identified the screening-off problem that any version must solve. \citet{Campbell2002} argued that naïve realism is needed to explain how perception makes demonstrative reference possible; \citet{Johnston2004} defended a variant on which hallucination presents uninstantiated property complexes; \citet{Fish2009} developed an eliminativist view of hallucinatory phenomenal character; \citet{Brewer2011} defended an object view on which illusion is explained by the conditions of perception. Critics argued that disjunctivism is unmotivated or that its account of the bad case is unstable \citep{Robinson1994,ByrneLogue2008}.

We take from this history the disjunctivist insight, that the good and bad cases need not share a fundamental kind, and supply what Martin's negative account of hallucination lacks: a positive characterization of the bad case.

\section{Framework}

\begin{definition}[Veridical perception; hallucination]\label{d:cases}
A \emph{veridical perception} is an experience in which a subject is presented with an object and some of its properties, as they are. A \emph{hallucination} is an experience as of an object and properties, occurring when no suitable object is perceived.
\end{definition}

\begin{definition}[Naïve realism]\label{d:nr}
\emph{Naïve realism} is the thesis that a veridical perceptual experience is a relation between subject and object of which the object and some of its properties are constituents, and that its phenomenal character is partly constituted by them.
\end{definition}

\begin{definition}[Fundamental kind]\label{d:kind}
The \emph{fundamental kind} of an experience is the most specific kind that says what the experience essentially is \citep{Martin2006}.
\end{definition}

\begin{definition}[Seeming]\label{d:seeming}
A \emph{seeming} is an attitude, with propositional content, in which it seems to the subject that things are a certain way. A seeming is not a belief, since one can have a seeming one disbelieves, as in the Müller-Lyer illusion.
\end{definition}

Definition~\ref{d:seeming} follows EPI-19, which holds that if it seems to a subject that $p$, she has prima facie justification for believing that $p$, and treats seemings as attitudes that perception and hallucination can produce alike. The present manuscript's thesis is about experiences; EPI-19's is about seemings, and the two are designed to fit together.

\begin{principle}[Presentational character]\label{p:pres}
Veridical perceptual experience is introspectively characterized as presenting its object, so that the object and its properties seem to the subject to be constituents of the experience and its character.
\end{principle}

\begin{principle}[Conservatism about the phenomenology]\label{p:cons}
A theory of perception that takes the introspective character of experience at face value is to be preferred, other things equal, to one that does not.
\end{principle}

Principle~\ref{p:pres} is the observation that \citet{Martin2002} develops from the transparency of experience: when one attends to one's perceptual experience, one finds the lemon and its yellowness, not an inner image. Principle~\ref{p:cons} is a principle of theory choice. It is consistent with MND-20, which holds that introspection is reliable about current phenomenal states.

\begin{principle}[Constitution of relational events]\label{p:events}
Whether an event is a relational event of a certain kind can depend on its relations to things beyond its local causes. Two events with the same local causes can therefore differ in fundamental kind.
\end{principle}

Principle~\ref{p:events} is general. A signature and a forged copy may involve the same movements of the hand; a payment and a mere transfer of coins, the same physical exchange; a handshake and a hand pressed against a dummy, the same muscular events in the person who extends the hand. What makes the event a handshake is in part that there is another hand to shake.

\section{The argument from hallucination}

The argument has two steps \citep{Robinson1994,Martin2004}.

\begin{enumerate}
\item[(A1)] For any veridical perception, there could be a hallucination indistinguishable from it by introspection, produced by the same proximal cause in the subject's brain.
\item[(A2)] \emph{Common kind}: if two experiences have the same proximal cause, they are of the same fundamental kind.
\item[(A3)] A hallucination does not have an external object as a constituent.
\item[(C)] Hence a veridical perception does not have an external object as a constituent, and naïve realism is false.
\end{enumerate}

(A1) is plausible on empirical grounds: stimulating the visual cortex produces visual experience, and the brain processes involved in perception are in principle reproducible in the absence of a stimulus. (A3) is true by Definition~\ref{d:cases}. The argument turns on (A2).

\begin{lemma}\label{l:ck}
The principle of common kind (A2) is false.
\end{lemma}

\begin{proof}
If veridical perception is a relational event, Principle~\ref{p:events} says that two events with the same local causes can differ in fundamental kind, depending on whether the relatum is present. (A2) therefore holds only if veridical perception is not a relational event; but whether it is relational is the point in dispute. So (A2) cannot be used without begging the question against naïve realism. Moreover, (A2) is false for relational events generally, as the examples given with Principle~\ref{p:events} show, and the defender of the argument must give a reason to think experiences an exception. The reason usually given is that the proximal cause determines the subject's experience because it determines how things seem to her. That reason supports a weaker principle, that the same proximal cause yields the same seeming, and we accept the weaker principle below. It does not support (A2).
\end{proof}

\begin{proposition}\label{p:nr}
Naïve realism is true.
\end{proposition}

\begin{proof}
By Principle~\ref{p:pres}, veridical experience is introspectively characterized as presenting its object as a constituent. Naïve realism takes this characterization at face value (Definition~\ref{d:nr}); its rivals do not. By Principle~\ref{p:cons}, naïve realism is to be preferred unless there is a decisive argument against it. The argument from hallucination is the main such argument, and by Lemma~\ref{l:ck} it relies on a false premise. The argument from illusion is considered in Section~\ref{s:illusion}. Hence naïve realism is true.
\end{proof}

\section{What hallucination is}

It follows from Proposition~\ref{p:nr} and (A3) that a hallucination is not of the same fundamental kind as a veridical perception. Martin's disjunctivism says nothing positive about the bad case: a hallucination is merely an experience indiscriminable by reflection from a veridical perception \citep{Martin2004,Martin2006}. This invites the \emph{screening-off} objection. If the brain state that produces a hallucination also occurs in the good case, and suffices for an experience indiscriminable from perception, then in the good case too there is such an experience, and it, rather than the relation to the object, seems to explain whatever the experience explains: the subject's beliefs, actions and introspective judgments. The relational experience is screened off.

We reply by characterizing hallucination positively and by locating the common factor elsewhere.

\begin{assumption}[Sensory imagery]\label{a:image}
A hallucination is a sensory-imagistic experience: an experience of the kind involved in vivid visual imagery and in dreams, produced in the absence of a suitable object by the processes that, in the good case, are engaged by the object. It has a phenomenal character of its own, which is intrinsic to it.
\end{assumption}

Assumption~\ref{a:image} is an empirical hypothesis supported by the overlap of neural mechanisms between visual imagery and visual perception and by the continuity of hallucination with imagery and dreaming. It fits MND-10, which holds that qualia are intrinsic to the experiences that have them: the hallucination has its qualia as an imagistic experience, and the veridical perception has its qualia with the perceived object as a constituent. MND-25 holds that dreams are conscious experiences; on Assumption~\ref{a:image} hallucinations are of the same broad kind as dreams.

\begin{lemma}[Seemings in common]\label{l:seem}
The same proximal cause produces the same seeming. In particular, a veridical perception of a lemon and a hallucination with the same proximal cause produce the same seeming that there is a yellow lemon there.
\end{lemma}

\begin{proof}
The content of a seeming is fixed by how the subject's cognitive system takes things to be, and this depends on the subject's states, not on whether a relatum outside them exists. Since the two experiences have the same proximal cause, and the subject's cognitive response to her experience is fixed by the proximal states it depends on, they produce the same seeming.
\end{proof}

\begin{theorem}\label{t:main}
Naïve realism is true, and hallucinations are experiences of a different kind from veridical perceptions. The two can nonetheless produce the same seeming, and their introspective indiscriminability is explained by that common seeming.
\end{theorem}

\begin{proof}
The first conjunct is Proposition~\ref{p:nr}. A veridical perception has its object as a constituent (Definition~\ref{d:nr}); a hallucination does not (A3), and is a sensory-imagistic experience (Assumption~\ref{a:image}); so their fundamental kinds differ. By Lemma~\ref{l:seem}, they can produce the same seeming. Introspective discrimination of two experiences proceeds by the subject's judging how things seem to her; if they seem the same, she cannot tell the experiences apart, even though they differ in kind.
\end{proof}

\begin{corollary}[Screening off]\label{c:screen}
The screening-off objection fails.
\end{corollary}

\begin{proof}
The objection requires that the proximal brain state in the good case suffice for a further experience, of the hallucinatory kind, alongside the perception. On Principle~\ref{p:events}, the same local state is, in the good case, part of the relational event that is the perception, and in the bad case part of an imagistic event; it does not constitute two experiences in the good case, any more than the movements of the hand constitute both a handshake and a separate failed handshake. What the good and bad cases have in common is the seeming of Lemma~\ref{l:seem}, which is not an experience but an attitude. The seeming explains what the two cases share: the subject's being inclined to believe that there is a lemon, her reaching for it, her inability to tell the cases apart. It does not screen off the relational experience from what only the good case explains: the subject's knowing that there is a lemon, her being able to refer demonstratively to that lemon \citep{Campbell2002}, and the phenomenal character that is constituted by the lemon.
\end{proof}

\section{Epistemological fit}

The distinction drawn in Theorem~\ref{t:main} between experiences, which differ in kind, and seemings, which can be common, fixes the relation between naïve realism and the epistemology of perception. EPI-19 holds that a seeming that $p$ gives prima facie justification for believing that $p$. By Lemma~\ref{l:seem}, the perceiver and the hallucinator have the same seeming and therefore the same prima facie justification. EPI-06 holds that justification is internalist while knowledge is externalist, and that an envatted subject is justified; this too is what Theorem~\ref{t:main} predicts, since justification goes with seemings and the envatted subject has them. The perceiver differs from the hallucinator in what she knows. EPI-21 holds that one's evidence is one's knowledge, so she also differs in her evidence; the difference is located in the relation to the object, not in a difference of seeming.

Naïve realism has sometimes been tied to the view that the perceiver's justification is better than the hallucinator's, because in the good case a fact is made manifest \citep{McDowell1982}. Theorem~\ref{t:main} does not require that view, and the fit with EPI-19 and EPI-06 is a reason to prefer a naïve realism that does without it. What the good case adds is knowledge and the evidence that knowledge constitutes, not a better seeming.

\section{Illusion and further cases}\label{s:illusion}

\paragraph{Illusion.} In the Müller-Lyer illusion, the lines are equal but one looks longer. The argument from illusion concludes that what one is aware of, the longer-looking line, is not the line itself. The naïve realist replies, following \citet{Brewer2011}, that the subject perceives the lines themselves, from a point of view and in conditions in which equal lines with outward fins have the visually relevant similarity to unequal lines. The experience is a relation to the lines; the seeming that one is longer is produced because of how the lines look from there. The illusion is then a case in which the experience is veridical in its object and the seeming is false. This is why one can disbelieve the seeming while continuing to see the lines.

\paragraph{Perspective.} Hume's diminishing table is perceived from a growing distance; its looking smaller is a relational property of the table, its occupying a smaller part of the visual field from a point of view, which the subject perceives. Nothing in the case requires an intermediary image.

\paragraph{Partial hallucination.} A subject who sees a real table on which a hallucinated cup appears to stand has an experience with two components: a perception of the table and an imagistic experience of a cup, integrated into one visual field. Theorem~\ref{t:main} applies to each component. MND-15 holds that consciousness is unified at each time in normal subjects, and partial hallucination is a case of one unified field with components of different kinds.

\paragraph{Veridical hallucination.} If an imagistic experience as of a lemon is produced by a brain stimulation while a lemon happens to be where it appears, the subject does not perceive the lemon, because the lemon plays no part in the event. On Definition~\ref{d:nr} the experience lacks the lemon as a constituent, and the verdict that she does not see it matches what the causal theory of perception requires \citep{Snowdon1981}.

\section{Recent literature}

\citet{ByrneLogue2008} distinguish metaphysical from epistemological disjunctivism and argue that the case for the former is weaker than often claimed; Theorem~\ref{t:main} is a version of metaphysical disjunctivism without the epistemological thesis. \citet{Logue2012} argues that naïve realism's best motivation is its explanation of phenomenal character; Principle~\ref{p:pres} is a version of that motivation. \citet{Soteriou2014} surveys the varieties of disjunctivism and the responses to the screening-off problem; our response belongs to the family that denies that the proximal state constitutes an experience in the good case, and is distinguished by locating the common factor in an attitude. \citet{Fish2009} holds that hallucinations lack phenomenal character altogether; Assumption~\ref{a:image}, which gives them an imagistic character of their own, is more conservative and more continuous with what is known about imagery.

\section{Objections and replies}

\begin{objection}[Indiscriminability requires identity]\label{o:ident}
If two experiences cannot be told apart by introspection, they have the same phenomenal character. Since naïve realism says the perception's character is constituted by its object, and the hallucination has no object, it must deny something introspection reveals.
\end{objection}

\begin{reply}
Indiscriminability does not entail sameness. It is not transitive, as series of color patches each indiscriminable from the next show, so it is not identity of anything \citep{Williamson1990}. Introspection reveals how things seem, and by Lemma~\ref{l:seem} the seemings are the same. It does not reveal whether the experience includes its object, which is not something introspection could reveal, since the difference is in a relation to something outside the subject.
\end{reply}

\begin{objection}[Causal argument]\label{o:causal}
The proximal cause in the brain must have an effect of the same kind in both cases, since the same cause yields the same effect \citep{Robinson1994}.
\end{objection}

\begin{reply}
The same cause yields the same \emph{local} effect, and we grant that the local effect is the same: the same neural events, and the same seeming. Whether that local effect is part of a perceptual relation depends on whether the object is suitably related to it (Principle~\ref{p:events}). The causal principle does not say that an effect's fundamental kind cannot depend on its relations to things other than its cause.
\end{reply}

\begin{objection}[Representationalism is simpler]\label{o:repr}
Representationalism gives one account of both cases: each experience represents a lemon. It is simpler than a disjunctive account.
\end{objection}

\begin{reply}
Representationalism gains simplicity by not taking Principle~\ref{p:pres} at face value: it says that the lemon merely seems to be a constituent of the experience. The disjunctive account is not more complex in what it posits, since representationalism also needs experiences, seemings and objects; it differs in recognizing that the perceptual relation constitutes the experience in the good case.
\end{reply}

\begin{objection}[Time lag]\label{o:lag}
Light from a distant star takes years to arrive, and the star may no longer exist. One cannot be related to what is not there, so perception of the star is not direct.
\end{objection}

\begin{reply}
Perception of the star is a relation to the star as it was when the light left it, and the relation obtains, since the star then existed and its existence is part of the causal chain. The case shows that perception can present its objects as they were, not that its objects are mental.
\end{reply}

\section{What the result does not establish}

The result does not establish that every perceptual experience is veridical, or that the perceiver has better justification than the hallucinator; it is compatible with their having the same justification, as EPI-19 and EPI-06 require. It does not decide which properties perception presents, such as whether natural kind properties are presented. It does not settle the details of the imagistic account of hallucination, which is an empirical hypothesis. The companion manuscript MND-16-b treats the argument from hallucination further.

\section{Conclusion}

Perception is direct. In veridical perception the object and its properties are constituents of the experience and its phenomenal character, as the experience itself presents them. The argument from hallucination fails because the principle of common kind is false for relational events. Hallucinations are experiences of a different kind, imagistic experiences with a character of their own, and what they share with perceptions is a seeming, an attitude in which things seem the same way. The distinction between experiences that differ in kind and seemings that can be common answers the screening-off objection and makes naïve realism consistent with a phenomenal conservative epistemology.

\begin{thebibliography}{99}\small
\bibitem[Ayer(1940)]{Ayer1940} Ayer, A.~J. (1940). \emph{The Foundations of Empirical Knowledge}. Macmillan.
\bibitem[Brewer(2011)]{Brewer2011} Brewer, B. (2011). \emph{Perception and Its Objects}. Oxford University Press.
\bibitem[Byrne and Logue(2008)]{ByrneLogue2008} Byrne, A. and Logue, H. (2008). Either/or. In A. Haddock and F. Macpherson (eds.), \emph{Disjunctivism: Perception, Action, Knowledge}, 57--94. Oxford University Press.
\bibitem[Campbell(2002)]{Campbell2002} Campbell, J. (2002). \emph{Reference and Consciousness}. Oxford University Press.
\bibitem[Fish(2009)]{Fish2009} Fish, W. (2009). \emph{Perception, Hallucination, and Illusion}. Oxford University Press.
\bibitem[Hinton(1973)]{Hinton1973} Hinton, J.~M. (1973). \emph{Experiences: An Inquiry into Some Ambiguities}. Clarendon Press.
\bibitem[Johnston(2004)]{Johnston2004} Johnston, M. (2004). The obscure object of hallucination. \emph{Philosophical Studies}, 120, 113--183.
\bibitem[Logue(2012)]{Logue2012} Logue, H. (2012). Why naive realism? \emph{Proceedings of the Aristotelian Society}, 112(2), 211--237.
\bibitem[Martin(2002)]{Martin2002} Martin, M.~G.~F. (2002). The transparency of experience. \emph{Mind and Language}, 17(4), 376--425.
\bibitem[Martin(2004)]{Martin2004} Martin, M.~G.~F. (2004). The limits of self-awareness. \emph{Philosophical Studies}, 120, 37--89.
\bibitem[Martin(2006)]{Martin2006} Martin, M.~G.~F. (2006). On being alienated. In T.~S. Gendler and J. Hawthorne (eds.), \emph{Perceptual Experience}, 354--410. Oxford University Press.
\bibitem[McDowell(1982)]{McDowell1982} McDowell, J. (1982). Criteria, defeasibility, and knowledge. \emph{Proceedings of the British Academy}, 68, 455--479.
\bibitem[Moore(1903)]{Moore1903} Moore, G.~E. (1903). The refutation of idealism. \emph{Mind}, 12(48), 433--453.
\bibitem[Robinson(1994)]{Robinson1994} Robinson, H. (1994). \emph{Perception}. Routledge.
\bibitem[Russell(1912)]{Russell1912} Russell, B. (1912). \emph{The Problems of Philosophy}. Williams and Norgate.
\bibitem[Siegel(2010)]{Siegel2010} Siegel, S. (2010). \emph{The Contents of Visual Experience}. Oxford University Press.
\bibitem[Snowdon(1981)]{Snowdon1981} Snowdon, P. (1981). Perception, vision and causation. \emph{Proceedings of the Aristotelian Society}, 81, 175--192.
\bibitem[Soteriou(2014)]{Soteriou2014} Soteriou, M. (2014). \emph{Disjunctivism}. Routledge.
\bibitem[Tye(1995)]{Tye1995} Tye, M. (1995). \emph{Ten Problems of Consciousness: A Representational Theory of the Phenomenal Mind}. MIT Press.
\bibitem[Williamson(1990)]{Williamson1990} Williamson, T. (1990). \emph{Identity and Discrimination}. Blackwell.
\end{thebibliography}

\end{document}
